% Page 065 — page imprimée 62
% Anecdotes, Témoignages, Souvenirs…

\begin{center}
{\bfseries\Large Anecdotes, Témoignages, Souvenirs…}
\end{center}

\medskip

\noindent{\bfseries EVA VERDIER} mère de Madeleine qui avait épousé Emilien Martin (originaire de
Rousses à 3km du Villaret dans la vallée du Bétuzon) raconte, à 104 ans, en 1986, à la
Bragousette où elle passa ses dernières années chez sa fille :

\medskip

\noindent{\bfseries La soie}

\begin{quote}
\textit{« Quand j’étais jeune j’habitais aux Salides, un petit hameau sur la pente sous le col Salidès.
Les mûriers étaient nombreux dans la région et on élevait des vers à soie dans beaucoup de
fermes. Les cocons étaient vendus à St. André de Valborgne, et je faisais éclore quelques
onces de « graine ». Dès que j’avais 5kg de cocons je les portais à la filature à Saint André. »}
\end{quote}

\begin{quote}
\textit{« La Fabrique » à Meyrueis était alors une grande magnanerie, puis elle fut transformée en
filature. »}
\end{quote}

\medskip

\noindent{\bfseries Un assassinat en 1826}

Celui de Monsieur de Fenouillet, Baron des Fons, (Les Fons est une grosse propriété à
quelques km. de Cabrillac sur la route de Saint André de Valborgne). C’est une histoire que
lui avait racontée son grand père :

\begin{quote}
\textit{« Monsieur de Fenouillet, propriétaire terrien vivant du revenu de ses terres cultivées par
des fermiers faisait confectionner par sa lingère habitant aux Salides (hameau en contrebas
du col Salidès) des sacs d’étoffe dans les quels il serrait les rouleaux d’écus qu’il amassait. Il
les cachait dans une pièce spéciale du château des Fons (était-ce son bureau ou sa cave, je ne
sais).}
\end{quote}

\begin{quote}
\textit{Trois jeunes gens des Salides décidèrent de le dévaliser. L’un d’eux, courtisant la lingère,
manœuvre habilement pour localiser le trésor, lui demandant d’abord de soustraire elle
même quelques sachets et se heurtant à un refus finit, après lui avoir demandé quel est son
travail principal, à quoi étaient destinés les sacs qu’elle cousait, comment et où ils étaient
rangés, à obtenir une nuit un rendez vous amoureux. Il la prie de laisser le soir une porte du
château ouverte.}
\end{quote}

\begin{quote}
\textit{Il pénétra alors la nuit avec ses deux acolytes et ils arrivèrent à la cachette, lorsque le baron
réveillé en sursaut par un bruit insolite se leva et leur barra la route. Pris de panique les trois
voleurs lui tranchèrent la gorge. Ils furent rapidement identifiés et condamnés à être
décapités.}
\end{quote}

\begin{quote}
\textit{La guillotine fut installée au Pompidou. Un jeune berger des Salides décida d’assister à
l’exécution pour « voir tomber les têtes ». Ce fut si horrible qu’il en perdit la raison. »}
\end{quote}
